\section{সমন্বিত সমস্যা ও নিজের দক্ষতা যাচাই}
প্রতিটি প্রশ্নে ভিত্তি, বিটের দৈর্ঘ্য এবং input/output-এর অর্থ আগে নির্ধারণ করো। পরে নিচের সমাধানের সঙ্গে নিজের ধাপ মিলিয়ে দেখো।

\chThreeExample{১}{মূল্য রূপান্তর ও যোগের মাধ্যমে বিয়োগ}
মিলি $(59)_{16}$ টাকা এবং রাফি $(106)_8$ টাকা দিয়ে বই কিনল। কে বেশি খরচ করেছে? আট বিটে ২-এর পরিপূরক দিয়ে ব্যবধান বের করো।
\chThreeAnswer মিলির খরচ $5\times16+9=89$; রাফির $1\times64+6=70$। আট বিটে $89=01011001$, $70=01000110$। ৭০-এর ২-এর পরিপূরক $10111010$।
\[01011001+10111010=1\;00010011.\]
ক্যারি বাদ দিয়ে $00010011=19$; মিলি ১৯ টাকা বেশি খরচ করেছে। দুই সংখ্যা ও ফল আট বিটের signed পরিসরের মধ্যে আছে।

\chThreeExample{২}{মিশ্র ভিত্তি, BCD এবং ক্যারেক্টার}
যন্ত্রে $(2D)_{16}$ দেখানো আছে। সেটিকে decimal, binary এবং BCD-তে লেখো।
\chThreeAnswer $2\times16+13=45$; binary 101101; BCD $0100\;0101$। display-এ লেখা ৪ ও ৫ দুটি ক্যারেক্টার; সংখ্যার binary pattern-এর সঙ্গে তাদের character code এক নয়।

\chThreeExample{৩}{কয়েকটি গেটের সমতুল্য সার্কিট}
$Y=AB+A\overline B+\overline AB$। সরল রাশি বের করে একটি gate-এ বাস্তবায়ন করো।
\chThreeAnswer $Y=A(B+\overline B)+\overline AB=A+\overline AB=A+B$। একটি OR যথেষ্ট। $A=B=0$-তে Y=০ এবং অন্য তিন সারিতে ১; মূল রাশির সঙ্গেও সব সারি মেলে।

\chThreeExample{৪}{কন্ট্রোল সার্কিটের গ অংশের মতো প্রয়োগ}
এক পরীক্ষাগারে মূল অনুমতি A থাকলে কাজ চলবে। অনুমতি না থাকলেও emergency switch B চালু হলে কাজ চলবে।
\chThreeAnswer মূল শর্ত $Y=A+\overline AB$। সরল করলে $Y=A+B$। A ও B-এর OR-এ output নেওয়া যায়। সমীকরণের ``অনুমতি নেই'' অংশটি আলাদা NOT-AND দিয়ে করলে একই সত্যক সারণি পাওয়া যায়, কিন্তু বেশি gate লাগে।

\chThreeExample{৫}{এনকোডার–ডিকোডার দিয়ে নির্বাচিত লাইন}
একটি আদর্শ one-hot 8-to-3 encoder-এ $D_5=1$। আউটপুট একটি enabled active-high 3-to-8 decoder-এ দেওয়া হলো।
\chThreeAnswer encoder-এর কোড 101; decoder-এর $Y_5=1$, অন্য line ০। এটি লাইন নির্বাচনের বিপরীত কাজের উদাহরণ; যেকোনো text/code encoding এবং display decoding একই তিন-bit circuit দিয়ে হয়—এ কথা বোঝায় না।

\chThreeExample{৬}{সময়মতো ডেটা ধরে রাখা}
এক sensor একবারে চার বিট 1011 তৈরি করে; পরের অংশ প্রতি ক্লকে একটি বিট নিতে পারে।
\chThreeAnswer PISO-তে এক ক্লকে 1011 load করে ধারাবাহিক SO পড়ো। ডান শিফটের সংযোগে এবং শিফটের আগে SO পড়লে বিট ১, ১, ০, ১। পরে আবার চার বিট একসঙ্গে দরকার হলে SIPO-তে একই order ও shift direction মেনে লোড করতে হবে।

\subsection{সংক্ষিপ্ত অনুধাবন যাচাই}
\begin{compactenum}
\item \textbf{বাইনারি কেন ডিজিটাল সার্কিটে সুবিধাজনক?} দুইটি পৃথক voltage range শনাক্ত করা সহজ। নির্ধারিত noise margin-এর মধ্যে সামান্য পরিবর্তন হলেও একই logic value শনাক্ত করা যায়; একেবারে noise-মুক্ত নয়।
\item \textbf{BCD কি সংখ্যা পদ্ধতি?} এটি প্রতিটি decimal digit-এর কোড; সাধারণ base-2 positional সংখ্যার পুরো মানের রূপান্তর নয়।
\item \textbf{একটি NOT কি universal gate?} NOT দিয়ে ইনপুট উল্টানো যায়, কিন্তু দুই ইনপুট AND/OR তৈরি করা যায় না। তাই একা universal নয়।
\item \textbf{Half adder-এর সীমা কী?} আগের ঘরের carry গ্রহণের তৃতীয় ইনপুট নেই। বহু বিটের ওপরের ঘরে full adder লাগে।
\item \textbf{Register ও counter কেন sequential?} তাদের বর্তমান stored state এবং নতুন input/control মিলে পরের state নির্ধারণ করে। শুধু বর্তমান input দিয়ে সব output নির্ধারিত হয় না।
\item \textbf{নাম Mod-10 হলেই synchronous কেন নয়?} ১০টি state-এর cycle modulus বোঝায়; synchronous/asynchronous নাম clock সংযোগ বোঝায়। দুটি শ্রেণিবিভাগ আলাদা।
\end{compactenum}
